El objetivo general de esta investigación es ampliar la comprensión sobre la planeación como función esencial de la administración y como elemento clave para orientar el desempeño de una organización. Se pretende analizar su concepto, su relevancia dentro del proceso administrativo, sus principios, sus etapas, sus tipos y las herramientas que permiten aplicarla de manera efectiva.

Entre los objetivos específicos podemos mencionar: definir el concepto de planeación y su relación con el proceso administrativo; explicar su importancia para anticipar problemas y orientar la toma de decisiones; describir sus principios y etapas; diferenciar los tipos de planeación según el nivel organizacional y el horizonte temporal; identificar herramientas útiles para su aplicación; y mostrar cómo la planeación contribuye a mejorar el desempeño de una empresa industrial.